\section{Background and related work}
\label{sec:background}

\subsection{Subnational microsimulation and spatial reweighting}

Subnational microsimulation usually starts from the same basic problem: a national survey contains
the behavioural and demographic detail needed for policy modelling, but it does not contain enough
observations in every local area to support direct estimation. Spatial microsimulation addresses
that gap either by reweighting existing survey records or by constructing synthetic small-area
populations from aggregate constraints \citep{tanton2014review, odonoghue2014review}.

This distinction matters for the present paper. Our pipeline remains a reweighting system built on
observed CPS households. We do clone households across sampled geographies, but the core empirical
problem is still calibration of survey-based microdata to administrative targets. That makes
classical calibration methods the closest reference point for the benchmark design.

The spatial microsimulation literature also provides broader context. \citet{williamson1998},
\citet{huang2001}, and \citet{harland2012} describe combinatorial and synthetic-population
approaches that search over record assignments area by area. \citet{tanton2011} and
\citet{lovelace2016} discuss reweighting approaches for small-area estimation and policy analysis.
These methods are useful precedents, but they usually target one geographic layer at a time. Our
setting requires one weighted microdata system to reproduce district-level, state-level, and
national targets jointly.

\subsection{Survey calibration}

Let $i = 1, \ldots, n$ index sampled units, let $d_i$ denote the initial design weight for unit
$i$, and let $\mathbf{x}_i \in \mathbb{R}^p$ denote the vector of auxiliary variables used for
calibration. Survey calibration seeks adjusted weights $w_i$ such that the weighted sample totals
match known population totals:
\begin{equation}
    \sum_{i=1}^{n} w_i \mathbf{x}_i = \mathbf{T},
    \label{eq:calibration_constraint}
\end{equation}
where $\mathbf{T} \in \mathbb{R}^p$ is the vector of published totals
\citep{deville1992, sarndal2007}. In practice, the auxiliary variables may be counts, indicators,
or continuous quantities such as income components. When the target system is internally
consistent, classical calibration methods seek exact reproduction of those totals.

The subnational problem in this paper extends that setup in two ways. First, the target vector
contains variables from multiple administrative sources and therefore mixes count and dollar
constraints. Second, those constraints are nested by geography: district-level totals should add to
state totals, which should add to national totals after uprating and reconciliation. In the full
production problem, that yields tens of thousands of targets with substantial collinearity across
rows.

\subsection{Generalized regression}

Generalized regression (GREG) calibration minimizes a distance from the initial weights while
enforcing the calibration equations \citep{deville1992, sarndal2007}. With the usual chi-squared
distance, the calibrated weight for unit $i$ can be written as
\begin{equation}
    w_i = d_i + d_i \mathbf{x}_i^\top
    \left(\sum_{j=1}^{n} d_j \mathbf{x}_j \mathbf{x}_j^\top\right)^{-1}
    \left(\mathbf{T} - \sum_{j=1}^{n} d_j \mathbf{x}_j\right).
    \label{eq:greg}
\end{equation}
GREG is attractive because it can accommodate quantitative targets directly and is well understood
in survey statistics.

Its weaknesses are also clear in a microsimulation setting. The matrix in
Equation~\ref{eq:greg} becomes harder to invert reliably as the target system grows and the rows
become more collinear. GREG can also assign negative weights. That is acceptable in some
estimation settings, but it is awkward in a microdata file that will later be used for tax-benefit
simulation, where analysts expect households to represent positive population mass.

\subsection{Iterative proportional fitting and raking}

Iterative proportional fitting \citep[IPF;][]{deming1940, ireland1968}, often called raking in
survey practice, adjusts weights multiplicatively so that selected margins match published totals.
When the margins are compatible and the algorithm converges, the fitted margins are matched
exactly. IPF preserves non-negativity and avoids matrix inversion, which makes it a natural
baseline for count-style calibration problems.

IPF is narrower than GREG in scope. Classical IPF is built for categorical or margin-style
constraints, not arbitrary linear dollar targets. It also does not naturally encode a large mixed
system of district, state, and national constraints in one pass. For that reason, the benchmark in
this paper uses IPF only on target subsets that fit its assumptions rather than forcing it onto the
full production matrix.

\subsection{Adjacent weighting literatures}

Two adjacent literatures inform the framing even though they are not the main empirical baselines
here. Small area estimation borrows strength across areas to improve local estimates from limited
samples \citep{anderson2013}. Covariate-balancing methods in causal inference choose weights to
balance observed covariates between treated and control units \citep{imai2014}. Both literatures
share the broader idea that weighting can repair deficiencies in the observed sample, but they
optimize for different inferential goals than the hierarchical administrative-target problem studied
here.

\subsection{\texorpdfstring{$L_0$}{L0} regularization and the Hard Concrete distribution}

$L_0$ regularization penalizes the number of active parameters in a model. In our setting, the
parameters are household-geography combinations rather than neural-network weights. Direct
optimization of the $L_0$ norm is combinatorial, so we use the Hard Concrete relaxation proposed by
\citet{louizos2018}. Each unit receives a gate $z_i \in [0, 1]$ that is sampled during training and
that collapses toward 0 or 1 after clipping. The expected number of active gates is differentiable,
which makes it possible to optimize calibration error and sparsity jointly.

That adaptation changes the role of regularization in calibration. Classical calibration methods
return one fitted weight vector for the chosen target system. The Hard Concrete gate makes the
number of retained records a direct object of optimization. This is the mechanism that lets one
pipeline produce both compact national datasets and larger subnational datasets.

\subsection{National-level predecessor}

\citet{woodruff2024} developed a national pipeline that imputes tax variables from the IRS Public
Use File onto CPS records and then calibrates the resulting dataset to national administrative
totals. The present paper extends that work in three directions. First, the pipeline now assigns
households to district-level units and states rather than keeping a single national geography.
Second, the optimizer now uses $L_0$ gates instead of dropout-style regularization, which produces
exact sparsity at inference time. Third, the paper evaluates the resulting system as a benchmark
study, with explicit comparisons to classical reference methods and an experimental design that
tracks what happens as the target system approaches production scale.
